% !TeX root = ../main.tex
%% =====  TEMATYKA I PROGRAM PRACY INŻYNIERSKIEJ ====

\newpage
\rightline{Kraków, 30 kwietnia 2026 r.}
\begin{center}
    {\bf Tematyka pracy inżynierskiej i~praktyki dyplomowej\\
        \makeatletter\@author\makeatother,
        studenta studiów pierwszego stopnia na kierunku teleinformatyka (nr albumu: 422842)}\\
\end{center}

\begin{center}
    Temat pracy inżynierskiej:\\
    \makeatletter
    {\bf \@title}\\
    \makeatother
\end{center}

\begin{tabular}{rl}
    Opiekun pracy:               & \supervisor \\
    Recenzent pracy:             & \reviewer \\
    Miejsce praktyki dyplomowej: & WIEiT AGH, Kraków \\
\end{tabular}

\begin{center}
    {\bf Program pracy inżynierskiej i~praktyki dyplomowej}
\end{center}

\begin{enumerate}
    \item Omówienie koncepcji i~zakresu pracy inżynierskiej z~opiekunem.
    \item Zebranie i~studium literatury z~zakresu pasywnej radiolokacji (PBR) oraz standardu transmisji Wi-Fi IEEE 802.11a (modulacja OFDM, estymacja kanału).
    \item Praktyka dyplomowa:
          \begin{itemize}
              \item zapoznanie się ze środowiskiem symulacyjnym MATLAB i~narzędziami przetwarzania sygnałów radarowych,
              \item implementacja modułów nadajnika Wi-Fi (ramka 802.11a Non-HT) oraz wielodrogowego kanału radiowego z~efektem Dopplera,
              \item opracowanie odbiornika pasywnego: synchronizacja Schmidl--Cox, demodulacja OFDM, estymacja kanału metodą Zero-Forcing, filtracja MTI oraz okienkowanie 2D Blackman--Harris.
          \end{itemize}
    \item Implementacja algorytmu koherentnego usuwania celów Coherent CLEAN oraz odbiornika korelacyjnego (funkcja wzajemnej niejednoznaczności CAF) jako metody odniesienia.
    \item Przeprowadzenie badań symulacyjnych efektywności detekcji radaru Wi-Fi: generacja map Range--Doppler (2D/3D), analiza przekrojów 1D, wyznaczenie charakterystyk $SNR_{\mathrm{out}} = f(SNR_{\mathrm{in}})$ oraz ocena odporności na zakłócenia i~wielodrogowość.
    \item Analiza porównawcza metody CIR i~CAF oraz ocena skuteczności algorytmu Coherent CLEAN w~scenariuszach wielocelowych.
    \item Opracowanie redakcyjne tekstu pracy inżynierskiej i~wniosków końcowych.
\end{enumerate}

\noindent
Termin oddania w~dziekanacie: \\[1cm]

\begin{center}
    \begin{tabular}{lcr}
        .............................................................. & ~~~ &
        ..............................................................                           \\
        (podpis kierownika katedry)                                    &     & (podpis opiekuna) \\
    \end{tabular}
\end{center}